% GENERATED FILE. Edit canonical lessons, then run npm run generate:books.
% canonical-source-hash: fnv1a64:69f4998ad5649e2d
% canonical-lessons: RU-C421-vzletat, RU-C421-vzlyot, RU-C421-prizemlyatsya, RU-C421-illyuminator, RU-C421-trap

\chapter{Take-off and landing}
\label{ch:ru-take-off-and-landing}
\hlchaptermodality{421}

\begin{chapteropening}
\textbf{By the end of this chapter:} \emph{I can talk about take-off and landing: taking off, the take-off, landing, the plane window and the steps.}

Complete the last lesson of chapter 421: I can talk about take-off and landing: taking off, the take-off, landing, the plane window and the steps.
\end{chapteropening}

\section[\texorpdfstring{\ru{взлетать} (vzletát)}{взлетать (vzletát)}]{\ru{взлетать} (vzletát) — to take off}

\label{lesson:RU-C421-vzletat}



\begin{quote}
Before the new one: say the Russian for a customs declaration, then the Russian for to declare.
\end{quote}

\subsection*{You'll want to know: \ru{взлетать}}
\textbf{\ru{взлетать}} (\emph{vzletát}) — \textquotedblleft{}to take off\textquotedblright{}.

\begin{quote}
\ru{Самолёт} \ru{медленно} \ru{взлетает} \ru{над} \ru{морем}.
\end{quote}

\begin{quote}
\emph{Samolyót médlenno vzletáyet nad mórem.}
\end{quote}

\begin{quote}
The plane slowly takes off over the sea.
\end{quote}

Imperfective; perfective \emph{\ru{взлететь}} (\emph{\ru{взлечу}, \ru{взлетишь}}). \emph{\ru{Вз}-} means up.

\begin{cousinweb}
\emph{\ru{взлетать}} is built from \emph{\ru{вз}-} 'up' and \emph{\ru{летать}} 'to fly'. EtymDB lists the noun \emph{\ru{взлёт}} from Proto-Slavic \emph{vŭzletŭ}.
\end{cousinweb}

\begin{grammarlens}[title={What you've built}]
One more word for telling, explaining and describing at B1.
\end{grammarlens}

\subsection*{Your turn}
\begin{itemize}
\raggedright
  \item \emph{Say it:} \emph{vzletát}
  \item \emph{Say it:} \emph{vzletát}, once more
  \item \emph{Say it:} a sentence of your own with \emph{vzletát}
  \item \emph{Recall:} say the Russian for a customs declaration, then the Russian for to declare, then say \emph{vzletát} again
\end{itemize}

\begin{tcolorbox}[breakable,colback=teal!4,colframe=teal!35!black,title={Before you move on}]
What does \ru{взлетать} mean? (\textquotedblleft{}To take off\textquotedblright{}.) Say it once more.
\end{tcolorbox}
\section[\texorpdfstring{\ru{взлёт} (vzlyot)}{взлёт (vzlyot)}]{\ru{взлёт} (vzlyot) — take-off}

\label{lesson:RU-C421-vzlyot}



\begin{quote}
Before the new one: say the Russian for to declare, then the Russian for to take off.
\end{quote}

\subsection*{You'll want to know: \ru{взлёт}}
\textbf{\ru{взлёт}} (\emph{vzlyot}) — \textquotedblleft{}take-off\textquotedblright{}.

\begin{quote}
\ru{Во} \ru{время} \ru{взлёта} \ru{нельзя} \ru{вставать}.
\end{quote}

\begin{quote}
\emph{Vo vrémya vzlyóta nelzyá vstavát.}
\end{quote}

\begin{quote}
You may not stand up during take-off.
\end{quote}

Masculine. \emph{\ru{Взлётная} \ru{полоса}} — the runway.

\begin{cousinweb}
\emph{\ru{взлёт}} is inherited from Proto-Slavic \emph{vŭzletŭ} (EtymDB), built from \emph{vŭz-} 'up' and the root of \emph{\ru{летать}} 'to fly'.
\end{cousinweb}

\begin{grammarlens}[title={What you've built}]
One more word for telling, explaining and describing at B1.
\end{grammarlens}

\subsection*{Your turn}
\begin{itemize}
\raggedright
  \item \emph{Say it:} \emph{vzlyot}
  \item \emph{Say it:} \emph{vzlyot}, once more
  \item \emph{Say it:} a sentence of your own with \emph{vzlyot}
  \item \emph{Recall:} say the Russian for to declare, then the Russian for to take off, then say \emph{vzlyot} again
\end{itemize}

\begin{tcolorbox}[breakable,colback=teal!4,colframe=teal!35!black,title={Before you move on}]
What does \ru{взлёт} mean? (\textquotedblleft{}Take-off\textquotedblright{}.) Say it once more.
\end{tcolorbox}
\section[\texorpdfstring{\ru{приземляться} (prizemlyátsya)}{приземляться (prizemlyátsya)}]{\ru{приземляться} (prizemlyátsya) — to land}

\label{lesson:RU-C421-prizemlyatsya}



\begin{quote}
Before the new one: say the Russian for to take off, then the Russian for take-off.
\end{quote}

\subsection*{You'll want to know: \ru{приземляться}}
\textbf{\ru{приземляться}} (\emph{prizemlyátsya}) — \textquotedblleft{}to land\textquotedblright{}.

\begin{quote}
\ru{Самолёт} \ru{приземляется} \ru{поздно} \ru{ночью}.
\end{quote}

\begin{quote}
\emph{Samolyót prizemlyáyetsya pózdno nóchyu.}
\end{quote}

\begin{quote}
The plane lands late at night.
\end{quote}

Imperfective; perfective \emph{\ru{приземлиться}}. Built on \emph{\ru{земля}} — the ground.

\begin{cousinweb}
\emph{\ru{приземляться}} is built from \emph{\ru{при}-} 'to', \emph{\ru{земля}} 'earth, the ground' and \emph{-\ru{ся}}. \emph{\ru{Земля}} is inherited from Proto-Slavic \emph{zemľa} 'earth, land' (EtymDB).
\end{cousinweb}

\begin{grammarlens}[title={What you've built}]
One more word for telling, explaining and describing at B1.
\end{grammarlens}

\subsection*{Your turn}
\begin{itemize}
\raggedright
  \item \emph{Say it:} \emph{prizemlyátsya}
  \item \emph{Say it:} \emph{prizemlyátsya}, once more
  \item \emph{Say it:} a sentence of your own with \emph{prizemlyátsya}
  \item \emph{Recall:} say the Russian for to take off, then the Russian for take-off, then say \emph{prizemlyátsya} again
\end{itemize}

\begin{tcolorbox}[breakable,colback=teal!4,colframe=teal!35!black,title={Before you move on}]
What does \ru{приземляться} mean? (\textquotedblleft{}To land\textquotedblright{}.) Say it once more.
\end{tcolorbox}
\section[\texorpdfstring{\ru{иллюминатор} (illyuminátor)}{иллюминатор (illyuminátor)}]{\ru{иллюминатор} (illyuminátor) — a porthole, a plane window}

\label{lesson:RU-C421-illyuminator}



\begin{quote}
Before the new one: say the Russian for take-off, then the Russian for to land.
\end{quote}

\subsection*{You'll want to know: \ru{иллюминатор}}
\textbf{\ru{иллюминатор}} (\emph{illyuminátor}) — \textquotedblleft{}a porthole, a plane window\textquotedblright{}.

\begin{quote}
\ru{Ребёнок} \ru{всю} \ru{дорогу} \ru{смотрел} \ru{в} \ru{иллюминатор}.
\end{quote}

\begin{quote}
\emph{Rebyónok vsyu dorógu smotrél v illyuminátor.}
\end{quote}

\begin{quote}
The child looked out of the window the whole way.
\end{quote}

Masculine. The window of a ship or a plane; a house window is \emph{\ru{окно}}.

\begin{cousinweb}
\emph{\ru{иллюминатор}} is usually traced to Latin \emph{illūmināre} 'to light up' (EtymDB): a round window that lets the light in. EtymDB does not list the Russian word, so its route is hedged.
\end{cousinweb}

\begin{grammarlens}[title={What you've built}]
One more word for telling, explaining and describing at B1.
\end{grammarlens}

\subsection*{Your turn}
\begin{itemize}
\raggedright
  \item \emph{Say it:} \emph{illyuminátor}
  \item \emph{Say it:} \emph{illyuminátor}, once more
  \item \emph{Say it:} a sentence of your own with \emph{illyuminátor}
  \item \emph{Recall:} say the Russian for take-off, then the Russian for to land, then say \emph{illyuminátor} again
\end{itemize}

\begin{tcolorbox}[breakable,colback=teal!4,colframe=teal!35!black,title={Before you move on}]
What does \ru{иллюминатор} mean? (\textquotedblleft{}A porthole, a plane window\textquotedblright{}.) Say it once more.
\end{tcolorbox}
\section[\texorpdfstring{\ru{трап} (trap)}{трап (trap)}]{\ru{трап} (trap) — the steps (to a plane), a gangway}

\label{lesson:RU-C421-trap}



\begin{quote}
Before the new one: say the Russian for to land, then the Russian for a porthole.
\end{quote}

\subsection*{You'll want to know: \ru{трап}}
\textbf{\ru{трап}} (\emph{trap}) — \textquotedblleft{}the steps (to a plane), a gangway\textquotedblright{}.

\begin{quote}
\ru{Пассажиры} \ru{спускаются} \ru{по} \ru{трапу}.
\end{quote}

\begin{quote}
\emph{Passazhíry spuskáyutsya po trápu.}
\end{quote}

\begin{quote}
The passengers come down the steps.
\end{quote}

Masculine. \emph{\ru{По} \ru{трапу}} — up or down the steps.

\begin{cousinweb}
\emph{\ru{трап}} was borrowed from Dutch \emph{trap} (EtymDB), the everyday Dutch word for stairs.
\end{cousinweb}

\begin{grammarlens}[title={What you've built}]
One more word for telling, explaining and describing at B1.
\end{grammarlens}

\subsection*{Your turn}
\begin{itemize}
\raggedright
  \item \emph{Say it:} \emph{trap}
  \item \emph{Say it:} \emph{trap}, once more
  \item \emph{Say it:} a sentence of your own with \emph{trap}
  \item \emph{Recall:} say the Russian for to land, then the Russian for a porthole, then say \emph{trap} again
\end{itemize}

\begin{tcolorbox}[breakable,colback=teal!4,colframe=teal!35!black,title={Before you move on}]
What does \ru{трап} mean? (\textquotedblleft{}The steps (to a plane), a gangway\textquotedblright{}.) Say it once more.
\end{tcolorbox}
